% GENERATED FILE. Edit canonical lessons, then run npm run generate:books.
% canonical-source-hash: fnv1a64:d169c2728820ea8a
% canonical-lessons: TA-C367-nanam, TA-C367-ariyamai, TA-C367-arinar, TA-C367-metai, TA-C367-poti

\chapter{Wisdom and Ignorance}
\label{ch:ta-wisdom-and-ignorance}
\hlchaptermodality{367}

\begin{chapteropening}
\textbf{By the end of this chapter:} \emph{I can talk about wisdom, ignorance and learned people.}

Complete the last lesson of chapter 367: I can talk about wisdom, ignorance and learned people.
\end{chapteropening}

\section[\texorpdfstring{\ta{ஞானம்} (ñāṉam)}{ஞானம் (ñāṉam)}]{\ta{ஞானம்} (ñāṉam) — wisdom}

\label{lesson:TA-C367-nanam}



\begin{quote}
Before the new one: say the Tamil for a formula, then the Tamil for to define.
\end{quote}

\subsection*{You'll want to know: \ta{ஞானம்}}
\textbf{\ta{ஞானம்}} — \emph{ñāṉam} — \textquotedblleft{}wisdom\textquotedblright{}.

\begin{quote}
\ta{அனுபவத்தால்} \ta{ஞானம்} \ta{வருகிறது}.
\end{quote}

\begin{quote}
\emph{aṉupavattāl ñāṉam varukiṟatu.}
\end{quote}

\begin{quote}
Wisdom comes with experience.
\end{quote}

A deeper word than \ta{அறிவு}, knowledge; \ta{ஞானி}, a sage.

\begin{cousinweb}
\textbf{\ta{ஞானம்}} (\emph{ñāṉam}) is borrowed from Sanskrit \emph{jñāna}, \textquotedblleft{}knowledge\textquotedblright{} (CDIAL 5281). Tamil writes the cluster jñ with one letter, \ta{ஞ}, as Pali does with \emph{ñāṇa}.
\end{cousinweb}

\begin{grammarlens}[title={What you've built}]
One more word for giving reasons and opinions.
\end{grammarlens}

\subsection*{Your turn}
\begin{itemize}
\raggedright
  \item \emph{Say it:} \emph{ñāṉam}
  \item \emph{Say it:} \emph{ñāṉam}, once more
  \item \emph{Say it:} \emph{aṉupavattāl ñāṉam varukiṟatu.} — Wisdom comes with experience.
  \item \emph{Recall:} say the Tamil for a formula, then the Tamil for to define, then say \emph{ñāṉam} again
\end{itemize}

\begin{tcolorbox}[breakable,colback=teal!4,colframe=teal!35!black,title={Before you move on}]
What does \ta{ஞானம்} mean? (\textquotedblleft{}Wisdom\textquotedblright{}.) Say it once more.
\end{tcolorbox}
\section[\texorpdfstring{\ta{அறியாமை} (aṟiyāmai)}{அறியாமை (aṟiyāmai)}]{\ta{அறியாமை} (aṟiyāmai) — ignorance}

\label{lesson:TA-C367-ariyamai}



\begin{quote}
Before the new one: say the Tamil for to define, then the Tamil for wisdom.
\end{quote}

\subsection*{You'll want to know: \ta{அறியாமை}}
\textbf{\ta{அறியாமை}} — \emph{aṟiyāmai} — \textquotedblleft{}ignorance\textquotedblright{}.

\begin{quote}
\ta{அவன்} \ta{அறியாமையால்} \ta{அந்தத்} \ta{தவறைச்} \ta{செய்தான்}.
\end{quote}

\begin{quote}
\emph{avaṉ aṟiyāmaiyāl antat tavaṟaic ceytāṉ.}
\end{quote}

\begin{quote}
He made that mistake out of ignorance.
\end{quote}

\ta{அறி}, to know + -\ta{ஆ}-, not + -\ta{மை}: not-knowing. \ta{பொறாமை}, envy, has the same shape.

\begin{cousinweb}
\textbf{\ta{அறியாமை}} (\emph{aṟiyāmai}) is native: \emph{aṟi}, \textquotedblleft{}to know\textquotedblright{} (DEDR 314), with the negative -\ta{ஆ}- and the abstract ending -\ta{மை}: \textquotedblleft{}the state of not knowing\textquotedblright{}.
\end{cousinweb}

\begin{grammarlens}[title={What you've built}]
One more word for giving reasons and opinions.
\end{grammarlens}

\subsection*{Your turn}
\begin{itemize}
\raggedright
  \item \emph{Say it:} \emph{aṟiyāmai}
  \item \emph{Say it:} \emph{aṟiyāmai}, once more
  \item \emph{Say it:} \emph{avaṉ aṟiyāmaiyāl antat tavaṟaic ceytāṉ.} — He made that mistake out of ignorance.
  \item \emph{Recall:} say the Tamil for to define, then the Tamil for wisdom, then say \emph{aṟiyāmai} again
\end{itemize}

\begin{tcolorbox}[breakable,colback=teal!4,colframe=teal!35!black,title={Before you move on}]
What does \ta{அறியாமை} mean? (\textquotedblleft{}Ignorance\textquotedblright{}.) Say it once more.
\end{tcolorbox}
\section[\texorpdfstring{\ta{அறிஞர்} (aṟiñar)}{அறிஞர் (aṟiñar)}]{\ta{அறிஞர்} (aṟiñar) — a scholar}

\label{lesson:TA-C367-arinar}



\begin{quote}
Before the new one: say the Tamil for wisdom, then the Tamil for ignorance.
\end{quote}

\subsection*{You'll want to know: \ta{அறிஞர்}}
\textbf{\ta{அறிஞர்}} — \emph{aṟiñar} — \textquotedblleft{}a scholar\textquotedblright{}.

\begin{quote}
\ta{அவர்} \ta{ஒரு} \ta{பெரிய} \ta{தமிழ்} \ta{அறிஞர்}.
\end{quote}

\begin{quote}
\emph{avar oru periya tamiḻ aṟiñar.}
\end{quote}

\begin{quote}
He is a great Tamil scholar.
\end{quote}

-\ta{ஞர்}, one who: \ta{கலைஞர்}, an artist, has the same ending.

\begin{cousinweb}
\textbf{\ta{அறிஞர்}} (\emph{aṟiñar}) is native: \emph{aṟi}, \textquotedblleft{}to know\textquotedblright{} (DEDR 314), with the ending -\ta{ஞர்}, \textquotedblleft{}one who\textquotedblright{}, as in \ta{கலைஞர்}: \textquotedblleft{}one who knows\textquotedblright{}.
\end{cousinweb}

\begin{grammarlens}[title={What you've built}]
One more word for giving reasons and opinions.
\end{grammarlens}

\subsection*{Your turn}
\begin{itemize}
\raggedright
  \item \emph{Say it:} \emph{aṟiñar}
  \item \emph{Say it:} \emph{aṟiñar}, once more
  \item \emph{Say it:} \emph{avar oru periya tamiḻ aṟiñar.} — He is a great Tamil scholar.
  \item \emph{Recall:} say the Tamil for wisdom, then the Tamil for ignorance, then say \emph{aṟiñar} again
\end{itemize}

\begin{tcolorbox}[breakable,colback=teal!4,colframe=teal!35!black,title={Before you move on}]
What does \ta{அறிஞர்} mean? (\textquotedblleft{}A scholar\textquotedblright{}.) Say it once more.
\end{tcolorbox}
\section[\texorpdfstring{\ta{மேதை} (mētai)}{மேதை (mētai)}]{\ta{மேதை} (mētai) — a genius}

\label{lesson:TA-C367-metai}



\begin{quote}
Before the new one: say the Tamil for ignorance, then the Tamil for a scholar.
\end{quote}

\subsection*{You'll want to know: \ta{மேதை}}
\textbf{\ta{மேதை}} — \emph{mētai} — \textquotedblleft{}a genius\textquotedblright{}.

\begin{quote}
\ta{அந்தச்} \ta{சிறுவன்} \ta{கணிதத்தில்} \ta{ஒரு} \ta{மேதை}.
\end{quote}

\begin{quote}
\emph{antac ciṟuvaṉ kaṇitattil oru mētai.}
\end{quote}

\begin{quote}
That boy is a genius at mathematics.
\end{quote}

X-\ta{இல்} \ta{மேதை}, a genius at X.

\begin{cousinweb}
\textbf{\ta{மேதை}} (\emph{mētai}) is borrowed from Sanskrit \emph{medhā}, \textquotedblleft{}mental vigour, intelligence, wisdom\textquotedblright{} (Monier-Williams); Tamil writes the final -ā as -\ta{ஐ} and the dh as \ta{த}.
\end{cousinweb}

\begin{grammarlens}[title={What you've built}]
One more word for giving reasons and opinions.
\end{grammarlens}

\subsection*{Your turn}
\begin{itemize}
\raggedright
  \item \emph{Say it:} \emph{mētai}
  \item \emph{Say it:} \emph{mētai}, once more
  \item \emph{Say it:} \emph{antac ciṟuvaṉ kaṇitattil oru mētai.} — That boy is a genius at mathematics.
  \item \emph{Recall:} say the Tamil for ignorance, then the Tamil for a scholar, then say \emph{mētai} again
\end{itemize}

\begin{tcolorbox}[breakable,colback=teal!4,colframe=teal!35!black,title={Before you move on}]
What does \ta{மேதை} mean? (\textquotedblleft{}A genius\textquotedblright{}.) Say it once more.
\end{tcolorbox}
\section[\texorpdfstring{\ta{போதி} (pōti)}{போதி (pōti)}]{\ta{போதி} (pōti) — to teach formally; to preach}

\label{lesson:TA-C367-poti}



\begin{quote}
Before the new one: say the Tamil for a scholar, then the Tamil for a genius.
\end{quote}

\subsection*{You'll want to know: \ta{போதி}}
\textbf{\ta{போதி}} — \emph{pōti} — \textquotedblleft{}to teach formally; to preach\textquotedblright{}.

\begin{quote}
\ta{அந்த} \ta{ஆசிரியர்} \ta{மாணவர்களுக்கு} \ta{நேர்மையைப்} \ta{போதித்தார்}.
\end{quote}

\begin{quote}
\emph{anta āciriyar māṇavarkaḷukku nērmaiyaip pōtittār.}
\end{quote}

\begin{quote}
That teacher taught the students honesty.
\end{quote}

Past (I …): \ta{போதித்தேன்}, \emph{pōtittēṉ}.

A formal, often moral teach: the person taught takes -\ta{க்கு}, the lesson -\ta{ஐ}.

\begin{cousinweb}
\textbf{\ta{போதி}} (\emph{pōti}) is a verb built on Sanskrit \emph{bodha}, \textquotedblleft{}waking, understanding\textquotedblright{} (CDIAL 9314), from the root \emph{budh}, \textquotedblleft{}to wake, to know\textquotedblright{}: to make someone understand.
\end{cousinweb}

\begin{grammarlens}[title={What you've built}]
One more word for giving reasons and opinions.
\end{grammarlens}

\subsection*{Your turn}
\begin{itemize}
\raggedright
  \item \emph{Say it:} \emph{pōti}
  \item \emph{Say it:} \emph{pōti}, once more
  \item \emph{Say it:} \emph{anta āciriyar māṇavarkaḷukku nērmaiyaip pōtittār.} — That teacher taught the students honesty.
  \item \emph{Recall:} say the Tamil for a scholar, then the Tamil for a genius, then say \emph{pōti} again
\end{itemize}

\begin{tcolorbox}[breakable,colback=teal!4,colframe=teal!35!black,title={Before you move on}]
What does \ta{போதி} mean? (\textquotedblleft{}To teach formally; to preach\textquotedblright{}.) Say it once more.
\end{tcolorbox}
